\section{Model Statute}
\label{sec:model}

The following model statute is drafted in U.S. Code style for purposes of federal enactment. State legislatures may adapt this text by substituting the relevant state agency for the Census Bureau, the relevant state redistricting body for the redistricting commission, and adding any state-specific adjustments required by the state constitution's substantive redistricting criteria.

\medskip

\noindent\rule{\linewidth}{0.4pt}

\bigskip

\noindent\textbf{THE DISTRICTING INTEGRITY ACT}

\textit{An Act to establish a uniform algorithmic method for congressional redistricting and to vest implementation authority in the Bureau of the Census.}

\noindent\textit{Be it enacted by the Senate and House of Representatives of the United States of America in Congress assembled,}

\bigskip

\noindent\textbf{SECTION 1. FINDINGS AND PURPOSE.}

\textbf{(a) Findings.} Congress finds that---

\begin{enumerate}[label=(\arabic*)]
    \item Congressional redistricting by partisan state legislatures has produced district maps designed to entrench political parties and insulate incumbents from electoral accountability, undermining the constitutional guarantee that the House of Representatives shall be chosen ``by the People.''

    \item The Supreme Court held in \textit{Rucho v.\ Common Cause}, 588 U.S. 684 (2019), that federal courts lack judicially manageable standards for adjudicating partisan gerrymandering claims, and identified federal legislation as an appropriate remedy.

    \item Congress has authority under Article I, Section 4 of the Constitution to regulate the manner of congressional elections, including the method by which congressional districts are drawn, as confirmed in \textit{Smiley v.\ Holm}, 285 U.S. 355 (1932).

    \item Mathematical algorithms that optimize for population equality and geographic compactness using only data from the decennial census are structurally incapable of incorporating partisan data and therefore structurally immune to partisan manipulation.

    \item The method of equal proportions (Huntington-Hill method), codified at section 2a of this title, demonstrates that Congress may successfully mandate mathematical methods for politically sensitive calculations, producing legitimate and stable outcomes that have operated without controversy since 1941.

    \item Algorithmic redistricting satisfies the requirements of section 2 of the Voting Rights Act of 1965 (52 U.S.C. \S\,10301) and produces on average more majority-minority districts than legislatively enacted plans.

    \item Public confidence in the fairness of congressional redistricting is a legitimate governmental interest that algorithmic redistricting advances by providing transparent, reproducible, and auditable district maps.
\end{enumerate}

\textbf{(b) Purpose.} The purpose of this Act is to establish an objective, transparent, and manipulation-resistant method for congressional redistricting that satisfies constitutional population-equality requirements, produces geographically compact and contiguous districts, and eliminates structural opportunities for partisan manipulation of district boundaries.

\bigskip

\noindent\textbf{SECTION 2. ALGORITHMIC REDISTRICTING METHOD.}

\textbf{(a) Definitions.} For purposes of this section---

\begin{enumerate}[label=(\arabic*)]
    \item The term ``certified redistricting algorithm'' means a computer program that---
    \begin{enumerate}[label=(\Alph*)]
        \item uses only data from the most recent decennial census conducted pursuant to section 141 of title 13 and geographic boundary data maintained by the Bureau under the Topologically Integrated Geographic Encoding and Referencing (TIGER) system;
        \item does not use, directly or indirectly, any data on the partisan affiliation, voting history, or electoral preferences of any person or geographic unit;
        \item produces congressional district maps that---
        \begin{enumerate}[label=(\roman*)]
            \item achieve equal population among districts, with maximum deviation not to exceed one-half of one percent of the average district population;
            \item ensure that all districts are geographically contiguous;
            \item minimize the ratio of district boundary perimeter to district area (maximize compactness) subject to the population-equality constraint; and
            \item do not divide census tracts except where necessary to satisfy population-equality requirements;
        \end{enumerate}
        \item produces identical output for identical input data; and
        \item has source code that is published and publicly accessible.
    \end{enumerate}

    \item The term ``preliminary district map'' means the congressional district map produced by the certified redistricting algorithm for a State without any human modification.

    \item The term ``final district map'' means the congressional district map adopted for a State under subsection (d).
\end{enumerate}

\textbf{(b) Certification of Algorithm.} Not later than 180 days after the date of enactment of this Act, the Bureau of the Census shall---

\begin{enumerate}[label=(\arabic*)]
    \item certify a redistricting algorithm that satisfies the requirements of subsection (a)(1);
    \item publish the certified algorithm's source code on a publicly accessible Federal website; and
    \item promulgate regulations establishing procedures for running the certified algorithm, reviewing its outputs, and requesting adjustments.
\end{enumerate}

\textbf{(c) Production of Preliminary District Maps.} Not later than 30 days after the Bureau transmits the apportionment statement required by section 2a(a) of this title, the Bureau shall---

\begin{enumerate}[label=(\arabic*)]
    \item run the certified redistricting algorithm for each State;
    \item produce a preliminary district map for each State;
    \item publish each preliminary district map and the underlying data and parameters on a publicly accessible Federal website; and
    \item transmit each preliminary district map to the Governor and presiding officer of each chamber of the legislature of the State for which the map was produced.
\end{enumerate}

\textbf{(d) Adjustment Period and Final Map Adoption.}

\begin{enumerate}[label=(\arabic*)]
    \item \textit{Comment period.} For a period of 60 days following the publication of the preliminary district map for a State, any person may submit to the Bureau proposed adjustments to the preliminary district map. Proposed adjustments must---
    \begin{enumerate}[label=(\Alph*)]
        \item be documented by reference to objective geographic, municipal, or community-of-interest factors;
        \item not introduce any data on the partisan affiliation, voting history, or electoral preferences of any person or geographic unit;
        \item maintain population equality within the limits specified in subsection (a)(1)(C)(i); and
        \item not reduce the compactness of any district more than 5 percent below the compactness of the corresponding district in the preliminary district map.
    \end{enumerate}

    \item \textit{Bureau review.} Not later than 30 days after the close of the comment period, the Bureau shall---
    \begin{enumerate}[label=(\Alph*)]
        \item review all proposed adjustments for compliance with paragraph (1);
        \item accept, modify, or reject each proposed adjustment, with written explanation; and
        \item publish the final district map for each State.
    \end{enumerate}

    \item \textit{Default rule.} If no timely adjustments are received, or if all adjustments are rejected, the preliminary district map shall serve as the final district map.
\end{enumerate}

\textbf{(e) Voting Rights Act Adjustment.} Notwithstanding subsection (b), a State may deviate from the certified redistricting algorithm's output to the extent necessary to comply with Section 2 of the Voting Rights Act of 1965 (52 U.S.C. \S\,10301), provided that---

\begin{enumerate}[label=(\arabic*)]
    \item the deviation is limited to the minimum necessary to achieve compliance with Section 2;
    \item the State publishes a written justification documenting satisfaction of the preconditions established in \textit{Thornburg v.\ Gingles}, 478 U.S. 30 (1986), for each minority community affected by the deviation;
    \item partisan affiliation data is not used in determining the adjustment, consistent with \textit{Louisiana v.\ Callais}, 608 U.S.\ \underline{\hspace{1em}} (2026); and
    \item the adjusted district map and the supporting Gingles justification are published on the same publicly accessible Federal website as the preliminary district map under subsection (c)(3) not later than 15 days before the final map adoption deadline under subsection (d)(2).
\end{enumerate}

\bigskip

\noindent\textbf{SECTION 3. JUDICIAL REVIEW AND ENFORCEMENT.}

\textbf{(a) Jurisdiction.} A three-judge district court convened pursuant to section 2284 of title 28 shall have original jurisdiction over any action challenging---

\begin{enumerate}[label=(\arabic*)]
    \item the certification of the redistricting algorithm under section 2(b);
    \item the preliminary or final district map produced for any State; or
    \item any adjustment accepted or rejected by the Bureau under section 2(d).
\end{enumerate}

\textbf{(b) Grounds for Challenge.} A challenge to a final district map shall be limited to claims that---

\begin{enumerate}[label=(\arabic*)]
    \item the certified algorithm failed to satisfy the requirements of section 2(a)(1);
    \item the Bureau failed to run the certified algorithm in accordance with the regulations promulgated under section 2(b)(3);
    \item the final district map fails to satisfy the population-equality, contiguity, or compactness requirements of section 2(a)(1)(C); or
    \item an accepted adjustment introduced data on partisan affiliation, voting history, or electoral preferences in violation of section 2(a)(1)(B).
\end{enumerate}

\textbf{(c) Time Limits.} Any action under this section shall be filed not later than 30 days after the publication of the final district map. The three-judge district court shall issue a final decision not later than 60 days after the filing of the complaint.

\textbf{(d) Appeals.} Appeal from a decision of the three-judge district court under this section shall lie directly to the Supreme Court pursuant to section 1253 of title 28.

\textbf{(e) Effective Date.} This Act shall take effect with respect to redistricting conducted following the first decennial census conducted after the date of enactment.

\noindent\rule{\linewidth}{0.4pt}

\bigskip

\subsection{Drafting Notes}

Several features of this model statute deserve explanation. First, the operative mandate is stated minimally: the statute specifies required properties of the algorithm (what it must use, what it must not use, what it must produce) without prescribing any particular algorithmic implementation.

\textit{Note on compactness measurement.} Section 2(a)(1)(C)(iii) defines compactness as minimizing ``the ratio of district boundary perimeter to district area''---the Polsby-Popper formulation. The reference implementation described in \citet{dellaLibera2026recursive} instead minimizes edge cut in the census-tract adjacency graph, which correlates with but is not identical to Polsby-Popper score. These two compactness criteria are related: districts with short borders in geographic space tend to cross fewer census-tract boundaries. For statutory purposes, we define ``maximum compactness'' as the minimum-edge-cut partition of the TIGER adjacency graph---an objective criterion independent of shapefile-specific perimeter measurement. Courts challenging a map on compactness grounds would evaluate whether the certified algorithm was run correctly on the published adjacency graph, not whether the Polsby-Popper score exceeds a specific threshold. State legislatures that prefer a Polsby-Popper standard may substitute ``minimize the ratio of district boundary perimeter to district area'' for ``minimize graph edge-cut'' in subsection (a)(1)(C)(iii), with the tradeoff that Polsby-Popper scores are sensitive to coastline and river-boundary complexity in ways that graph edge-cut is not.

The overall property-specification approach is intentional---it allows the Bureau to certify the best available algorithm at the time of implementation and to recertify improved algorithms as computer science advances, without requiring statutory amendment. The analogy is the Bureau's authority to conduct the census using updated statistical methodology; the statute specifies the census's purpose and output requirements, not its technical implementation.

Second, the adjustment provision in Section 2(d) is deliberately narrow. It permits adjustments based on documented geographic and community factors but explicitly prohibits the introduction of partisan data at the adjustment stage. This prevents the adjustment process from becoming a backdoor for the partisan manipulation the statute is designed to eliminate. The 5\% compactness floor ensures that adjustments cannot substantially degrade the algorithm's outputs.

Third, the judicial review provision in Section 3 is administratively manageable---a deliberate contrast to the partisan gerrymandering standard that \textit{Rucho} found unmanageable. Courts review whether the algorithm was certified correctly, whether it was run correctly, and whether the outputs meet the statutory technical specifications. These are questions courts can answer without engaging in the normative balancing that defeated judicial oversight of partisan gerrymandering.
